% Generated by publication/build_sections.py; do not edit by hand.
\paragraph{Informal verdict.}
\textbf{A --- Complete solution}, with high confidence.

\paragraph{Lean coverage.}
Lean proves the complete canonical Q759 statement: for every fixed derivative order, both displayed sequences converge uniformly on [-1,1] to the corresponding derivative of the printed alpha. The response-only quantitative rates and sharpness claim are not formalized because the PDF does not request them.

\subsection{Audited informal answer}
\subsubsection{Faithful restatement}

Let
\[
E=C^\infty([-1,1];\mathbb R),\qquad
\lVert h\rVert_\infty=\max_{x\in[-1,1]}|h(x)|.
\]
For a sequence \((f_n)\subset E\) and \(g\in E\), the notation
\[
(f_n)\Rightarrow g
\]
means
\[
\forall k\in\mathbb N,\qquad
\lim_{n\to\infty}|f_n^{(k)}-g^{(k)}|_\infty=0.
\]

Define
\[
\alpha(x)=
\begin{cases}
\displaystyle \exp\!\left(-\frac{x^2}{1-x^2}\right),& |x|<1,\\[1ex]
0,&x=\pm1.
\end{cases}
\]
Determine whether \(f_n\Rightarrow\alpha\) for each of the two sequences
\[
f_n^{(a)}(x)=\exp\!\left(-x^2-x^4-\cdots-x^{2n}\right)
\]
and
\[
f_n^{(b)}(x)=
\left(1-\frac{x^2+x^4+\cdots+x^{2n}}{n}\right)^n.
\]

The visual PDF confirms that in part \(b\) the whole sum \(x^2+x^4+\cdots+x^{2n}\) is divided by \(n\). There is no apparent mathematical typo. The phrase ``equip \(E\) with the uniform norm'' is to be read together with the displayed definition: \(\Rightarrow\) is convergence for all seminorms \(h\mapsto|h^{(k)}|_\infty\), i.e. the usual \(C^\infty\)-topology. 

\medskip\hrule\medskip

\subsubsection{Final result}

\paragraph{Theorem}

Both sequences converge to \(\alpha\) in the sense of the problem:
\[
\boxed{f_n^{(a)}\Rightarrow\alpha}
\qquad\text{and}\qquad
\boxed{f_n^{(b)}\Rightarrow\alpha}.
\]

More precisely, for every \(k,M\in\mathbb N\),
\[
\max_{0\le j\le k}
\left\lvert \bigl(f_n^{(a)}-\alpha\bigr)^{(j)}\right\rVert_\infty
=O_{k,M}(n^{-M}),
\]
whereas, for every fixed \(k\),
\[
\max_{0\le j\le k}
\left\lvert \bigl(f_n^{(b)}-\alpha\bigr)^{(j)}\right\rvert _\infty
=O_k(n^{-1}).
\]
The rate \(O(n^{-1})\) in the second assertion is sharp already for uniform convergence of the functions.

\medskip\hrule\medskip

\subsection{1. Reduction to functions on ([0,1])}

Set
\[
P_n(t)=\sum_{j=1}^n t^j,\qquad 0\le t\le1,
\]
and, for \(0\le t<1\),
\[
P(t)=\sum_{j=1}^\infty t^j=\frac{t}{1-t}.
\]
Define
\[
A_n(t)=e^{-P_n(t)},\qquad
B_n(t)=\left(1-\frac{P_n(t)}n\right)^n,
\]
and
\[
\beta(t)=
\begin{cases}
e^{-P(t)}=\displaystyle e^{-t/(1-t)},&0\le t<1,\\
0,&t=1.
\end{cases}
\]
Then
\[
f_n^{(a)}(x)=A_n(x^2),\qquad
f_n^{(b)}(x)=B_n(x^2),\qquad
\alpha(x)=\beta(x^2).
\]

Thus it suffices to prove
\[
A_n\longrightarrow\beta
\quad\text{and}\quad
B_n\longrightarrow\beta
\qquad\text{in }C^\infty([0,1]).
\]

Indeed, for every smooth \(U\) and every \(k\ge0\), repeated differentiation gives
\[
\frac{d^k}{dx^k}U(x^2)
=
%
k!\sum_{\ell=0}^{\lfloor k/2\rfloor}
\frac{(2x)^{k-2\ell}}
{(k-2\ell)!,\ell\!},
U^{(k-\ell)}(x^2).
\]
Consequently,
\[
\left\lvert \frac{d^k}{dx^k}\bigl(U_n(x^2)-U(x^2)\bigr)\right\rVert_\infty
\le C_k\max_{0\le j\le k}|U_n^{(j)}-U^{(j)}|_{\infty,[0,1]}.
\]

\medskip\hrule\medskip

\subsubsection{1.1. Smoothness and flatness of \(\beta\)}

For \(t<1\), put
\[
u=\frac1{1-t}.
\]
Then
\[
\beta(t)=e^{1-u}=e,e^{-u}.
\]
Since \(u'=u^2\), induction gives polynomials \(Q_r\) such that
\[
\beta^{(r)}(t)=\beta(t),Q_r!\left(\frac1{1-t}\right),
\qquad 0\le t<1.
\]
Indeed, \(Q_0=1\), and if the formula holds at order \(r\), then
\[
Q_{r+1}(u)=u^2\bigl(Q_r'(u)-Q_r(u)\bigr).
\]

For every polynomial \(Q\),
\[
e^{-u}Q(u)\longrightarrow0\qquad(u\to+\infty).
\]
Hence
\[
\lim_{t\to1^-}\beta^{(r)}(t)=0
\qquad\text{for every }r\ge0.
\]
Extending every derivative by the value \(0\) at \(t=1\) gives a \(C^\infty\) function. For example, differentiability at the endpoint follows successively from the mean value theorem. Thus
\[
\beta\in C^\infty([0,1]),\qquad
\beta^{(r)}(1)=0\quad(r\ge0).
\]
In particular, \(\alpha=\beta\circ(x\mapsto x^2)\) belongs to \(E\) and is flat at both endpoints.

\medskip\hrule\medskip

\subsection{2. The first sequence}

We prove
\[
A_n=e^{-P_n}\longrightarrow\beta=e^{-P}
\quad\text{in }C^\infty([0,1]).
\]

\subsubsection{2.1. Elementary derivative bounds}

For \(r\ge0\) and \(0\le t\le1\),
\[
|P_n^{(r)}(t)|\le n^{r+1}.
\]
For \(r=0\) this is \(P_n(t)\le n\). For \(r\ge1\),
\[
P_n^{(r)}(t)
=\sum_{j=r}^n (j)_r t^{j-r},
\]
where
\[
(j)_r=j(j-1)\cdots(j-r+1),
\]
and hence
\[
|P_n^{(r)}(t)|
\le\sum_{j=1}^n j^r
\le n^{r+1}.
\]

Repeated use of the product and chain rules shows that every term in \(A_n^{(r)}\) is of the form
\[
c,e^{-P_n}
\prod_{i=1}^{m}P_n^{(s_i)},
\]
where
\[
s_i\ge1,\qquad
s_1+\cdots+s_m=r,\qquad
m\le r.
\]
Therefore
\[
\boxed{
|A_n^{(r)}(t)|
\le C_r n^{2r}e^{-P_n(t)}
\le C_r n^{2r}.}
\tag{2.1}
\]

\medskip\hrule\medskip

\subsubsection{2.2. Boundary layer}

Let
\[
\rho_n=n^{-1/2},\qquad
J_n=[1-\rho_n,1].
\]

We first show that
\[
P_n(t)\ge c\sqrt n
\qquad(t\in J_n)
\tag{2.2}
\]
for an absolute \(c>0\) and all sufficiently large \(n\).

Choose
\[
m=\left\lfloor\frac1{4\rho_n}\right\rfloor.
\]
Then \(m\le n\) for large \(n\). If \(t\ge1-\rho_n\) and \(1\le j\le m\), then, using
\[
\log(1-s)\ge-2s\qquad(0\le s\le\tfrac12),
\]
we get
\[
t^j\ge(1-\rho_n)^m
\ge e^{-2m\rho_n}
\ge e^{-1/2}.
\]
Moreover \(m\ge(8\rho_n)^{-1}\) for large \(n\). Thus
\[
P_n(t)\ge\sum_{j=1}^m t^j
\ge \frac{e^{-1/2}}{8\rho_n}
=c\sqrt n.
\]

Combining (2.1) and (2.2),
\[
\sup_{t\in J_n}|A_n^{(r)}(t)|
\le C_r n^{2r}e^{-c\sqrt n}.
\tag{2.3}
\]

For \(\beta\), the representation
\[
\beta^{(r)}(t)
=\beta(t)Q_r!\left(\frac1{1-t}\right)
\]
and \(1/(1-t)\ge\sqrt n\) on \(J_n\) imply
\[
\sup_{t\in J_n}|\beta^{(r)}(t)|
\le C_r e^{-c_r\sqrt n}.
\tag{2.4}
\]

Thus all derivatives of both \(A_n\) and \(\beta\) are uniformly negligible in the boundary layer.

\medskip\hrule\medskip

\subsubsection{2.3. Interior region}

Let
\[
I_n=[0,1-\rho_n].
\]
For \(t<1\),
\[
P(t)-P_n(t)
=\sum_{j=n+1}^\infty t^j
=\frac{t^{n+1}}{1-t}
=:R_n(t).
\]
Hence
\[
\beta(t)=A_n(t)e^{-R_n(t)}
\]
and
\[
A_n(t)-\beta(t)
=A_n(t)\bigl(1-e^{-R_n(t)}\bigr).
\tag{2.5}
\]

For \(r\ge0\), Leibniz's rule applied to
\[
R_n(t)=t^{n+1}(1-t)^{-1}
\]
gives
\[
R_n^{(r)}(t)
=
%
\sum_{a=0}^r
\binom ra (n+1)_a t^{n+1-a}
(r-a)!(1-t)^{-(r-a+1)}.
\]
On \(I_n\), \(1-t\ge\rho_n\) and \(t\le1-\rho_n\). Therefore
\[
|R_n^{(r)}(t)|
\le
C_r n^r\rho_n^{-(r+1)}
(1-\rho_n)^{n+1-r}.
\]
Since
\[
(1-\rho_n)^{n+1-r}
\le e^{-\rho_n(n+1-r)}
\le e^{-\sqrt n/2}
\]
for large \(n\), we obtain
\[
\max_{0\le r\le k}
|R_n^{(r)}|_{\infty,I_n}
\le C_k n^{C_k}e^{-\sqrt n/2}.
\tag{2.6}
\]

Every derivative of \(1-e^{-R_n}\) is a finite sum of products of derivatives of \(R_n\), with at least one such factor. Thus (2.6) implies
\[
\max_{0\le r\le k}
\left\lvert \bigl(1-e^{-R_n}\bigr)^{(r)}\right\rvert _{\infty,I_n}
\le C_k n^{C_k}e^{-c\sqrt n}.
\tag{2.7}
\]

Applying Leibniz's rule to (2.5), using (2.1) and (2.7), gives
\[
\max_{0\le r\le k}
|A_n^{(r)}-\beta^{(r)}|_{\infty,I_n}
\le C_k n^{C_k}e^{-c_k\sqrt n}.
\tag{2.8}
\]

Together with the boundary estimates (2.3)--(2.4), this proves
\[
\boxed{
|A_n-\beta|_{C^k([0,1])}
\le C_k n^{C_k}e^{-c_k\sqrt n}.}
\tag{2.9}
\]
In particular, for every \(M\),
\[
|A_n-\beta|_{C^k([0,1])}=O_{k,M}(n^{-M}).
\]

After composition with \(x^2\), this proves part \(a\):
\[
\boxed{f_n^{(a)}\Rightarrow\alpha.}
\]

\medskip\hrule\medskip

\subsection{3. A weighted binomial approximation}

To handle the second sequence uniformly up to the endpoint, we need a version of
\[
\left(1-\frac yn\right)^n\longrightarrow e^{-y}
\]
that remains valid after differentiation and permits \(y\) to range over ([0,n]).

\subsubsection{Lemma 3.1}

For every pair of nonnegative integers (m,L), there is a constant \(C_{m,L}\) such that, for all sufficiently large \(n\),
\[
\boxed{
\sup_{0\le y\le n}
(1+y)^L
\left\lvert 
\frac{d^m}{dy^m}\left(1-\frac yn\right)^n
-\frac{d^m}{dy^m}e^{-y}
\right\rvert 
\le\frac{C_{m,L}}n.}
\tag{3.1}
\]

\paragraph{Proof}

For \(n\ge m\),
\[
\frac{d^m}{dy^m}\left(1-\frac yn\right)^n
=(-1)^m c_{n,m}\left(1-\frac yn\right)^{n-m},
\]
where
\[
c_{n,m}=\frac{(n)_m}{n^m}
=\prod_{j=0}^{m-1}\left(1-\frac jn\right).
\]
Also
\[
\frac{d^m}{dy^m}e^{-y}=(-1)^m e^{-y}.
\]
Since all factors in \(c_{n,m}\) lie in ([0,1]),
\[
0\le1-c_{n,m}
\le\sum_{j=0}^{m-1}\frac jn
\le\frac{C_m}{n}.
\tag{3.2}
\]

Put
\[
h_{n,m}(y)=\left(1-\frac yn\right)^{n-m}.
\]

\paragraph{Region \(0\le y\le n/2\)}

Set \(z=y/n\) and
\[
a=-(n-m)\log(1-z),
\]
so that \(h_{n,m}(y)=e^{-a}\). For \(n\ge2m\),
\[
a\ge(n-m)z\ge\frac y2.
\]
Moreover,
\[
a-y
=(n-m)\bigl(-\log(1-z)-z\bigr)-mz.
\]
For \(0\le z\le1/2\),
\[
0\le-\log(1-z)-z
=\sum_{\ell=2}^\infty\frac{z^\ell}{\ell}
\le\frac{z^2}{1-z}
\le2z^2.
\]
Consequently,
\[
|a-y|
\le C_m\left(\frac yn+\frac{y^2}{n}\right).
\]
By the mean value theorem applied to \(s\mapsto e^{-s}\),
\[
|e^{-a}-e^{-y}|
\le e^{-\min(a,y)}|a-y|
\le
\frac{C_m}{n}(y+y^2)e^{-y/2}.
\]
Multiplication by \((1+y)^L\) preserves the \(O(1/n)\) bound, because every polynomial times \(e^{-y/2}\) is bounded on \([0,\infty\)).

Also
\[
h_{n,m}(y)\le e^{-y/2},
\]
so (3.2) gives the same weighted \(O(1/n)\) estimate for
\[
(c_{n,m}-1)h_{n,m}(y).
\]

\paragraph{Region \(n/2\le y\le n\)}

Here
\[
h_{n,m}(y)\le2^{-(n-m)},\qquad
e^{-y}\le e^{-n/2}.
\]
Thus
\[
(1+y)^L\bigl(h_{n,m}(y)+e^{-y}\bigr)
\le
(1+n)^L\bigl(2^{-(n-m)}+e^{-n/2}\bigr)
=O_{m,L}(n^{-1}).
\]

Combining both regions and (3.2) proves the lemma. \(\square\)

\medskip\hrule\medskip

\subsection{4. Derivatives of the geometric partial sums}

The following estimate allows us to compose Lemma 3.1 with \(P_n\).

\subsubsection{Lemma 4.1}

For every \(r\ge0\), there is \(C_r>0\), independent of \(n\), such that
\[
\boxed{
|P_n^{(r)}(t)|
\le C_r\bigl(1+P_n(t)\bigr)^{r+1}}
\tag{4.1}
\]
for all \(n\ge1\) and \(0\le t\le1\).

\paragraph{Proof}

The case \(r=0\) is immediate. Suppose \(r\ge1\).

If \(0\le t\le1/2\), then
\[
P_n^{(r)}(t)
\le\sum_{j=r}^\infty (j)_r2^{-(j-r)}
=:C_r,
\]
and (4.1) follows.

Now suppose \(t>1/2\), and set
\[
\delta=1-t,\qquad
L=\min(n,\delta^{-1}).
\]

First,
\[
P_n^{(r)}(t)
\le\sum_{j=1}^n j^r
\le C_r n^{r+1}.
\tag{4.2}
\]
Also, since \(t^{-r}\le2^r\) and \(t^j\le e^{-\delta j}\),
\[
P_n^{(r)}(t)
\le
2^r\sum_{j=1}^\infty j^r e^{-\delta j}
\le C_r\delta^{-(r+1)}.
\tag{4.3}
\]
Therefore
\[
P_n^{(r)}(t)\le C_r L^{r+1}.
\tag{4.4}
\]

It remains to compare \(L\) to \(P_n(t)\). If \(L\le8\), then trivially
\[
L\le8(1+P_n(t)).
\]
If \(L>8\), choose
\[
m=\left\lfloor\frac L4\right\rfloor.
\]
Then
\[
m\ge\frac L8,\qquad m\le n,\qquad m\delta\le\frac14.
\]
For \(1\le j\le m\),
\[
t^j=(1-\delta)^j
\ge e^{-2\delta j}
\ge e^{-1/2}.
\]
Hence
\[
P_n(t)\ge\sum_{j=1}^m t^j
\ge e^{-1/2}m
\ge cL.
\]
Combining this with (4.4) proves (4.1). \(\square\)

\medskip\hrule\medskip

\subsection{5. The second sequence}

Let
\[
G_n(y)=\left(1-\frac yn\right)^n,\qquad 0\le y\le n,
\]
and
\[
G(y)=e^{-y}.
\]
Then
\[
B_n=G_n\circ P_n,\qquad
A_n=G\circ P_n.
\]

We prove that for every fixed \(k\),
\[
|B_n-A_n|_{C^k([0,1])}=O_k(n^{-1}).
\tag{5.1}
\]

For \(k=0\), Lemma 3.1 with \(m=L=0\) immediately gives
\[
|B_n-A_n|_\infty\le\frac Cn.
\]

Now let \(k\ge1\). Repeated application of the chain rule shows that the \(k\)-th derivative of
\[
(G_n-G)\circ P_n
\]
is a finite linear combination, with coefficients depending only on \(k\), of terms
\[
\bigl(G_n^{(m)}-G^{(m)}\bigr)(P_n(t))
\prod_{i=1}^{m}P_n^{(r_i)}(t),
\]
where
\[
1\le m\le k,\qquad
r_i\ge1,\qquad
r_1+\cdots+r_m=k.
\]

By Lemma 4.1,
\[
\prod_{i=1}^m |P_n^{(r_i)}(t)|
\le
C_k
\bigl(1+P_n(t)\bigr)^{\sum_i(r_i+1)}
=
%
C_k
\bigl(1+P_n(t)\bigr)^{k+m}.
\]
Because \(m\le k\),
\[
\prod_{i=1}^m |P_n^{(r_i)}(t)|
\le
C_k\bigl(1+P_n(t)\bigr)^{2k}.
\tag{5.2}
\]

Applying Lemma 3.1 with weight exponent \(L=2k\) gives
\[
\left\lvert 
G_n^{(m)}(P_n(t))-G^{(m)}(P_n(t))
\right\rvert 
\le
\frac{C_k}
{n\bigl(1+P_n(t)\bigr)^{2k}}.
\tag{5.3}
\]
The product of (5.2) and (5.3) is \(O_k(1/n)\), uniformly in \(t\). There are only finitely many chain-rule terms, so (5.1) follows.

From the first part,
\[
|A_n-\beta|_{C^k([0,1])}=O_{k,M}(n^{-M})
\]
for every \(M\). Taking \(M=2\), for example,
\[
\boxed{
|B_n-\beta|_{C^k([0,1])}
\le\frac{C_k}{n}.}
\tag{5.4}
\]

Composing with \(x^2\) now yields
\[
\boxed{f_n^{(b)}\Rightarrow\alpha.}
\]

\medskip\hrule\medskip

\subsection{6. Boundary checks, small cases, and sharpness}

\subsubsection{6.1. At \(x=0\)}

For every \(n\),
\[
f_n^{(a)}(0)=f_n^{(b)}(0)=\alpha(0)=1.
\]
All three functions are even, so all odd derivatives at \(0\) vanish.

For the first sequence, the agreement is much stronger. Near \(0\),
\[
\sum_{j=n+1}^\infty x^{2j}=O(x^{2n+2}),
\]
and therefore
\[
f_n^{(a)}(x)-\alpha(x)=O(x^{2n+2}).
\]
Thus the derivatives at \(0\) agree exactly through order \(2n+1\).

\medskip\hrule\medskip

\subsubsection{6.2. At \(x=\pm1\)}

Since \(\alpha\) is flat,
\[
\alpha^{(k)}(\pm1)=0
\qquad(k\ge0).
\]

For the first sequence,
\[
f_n^{(a)}(\pm1)=e^{-n}.
\]
Every fixed-order derivative at \(\pm1\) is \(e^{-n}\) times a polynomial in \(n\), and hence tends to \(0\).

For the second sequence,
\[
1-\frac{x^2+x^4+\cdots+x^{2n}}n
\]
has a simple zero at each of \(x=\pm1\). Hence \(f_n^{(b)}\) has a zero of order \(n\) at each endpoint. In particular,
\[
\bigl(f_n^{(b)}\bigr)^{(k)}(\pm1)=0
\qquad\text{whenever }n>k.
\]

\medskip\hrule\medskip

\subsubsection{6.3. Sharpness for part \(b\)}

Fix \(0<t<1\). Since \(P_n(t)\to P(t)\) exponentially fast and remains bounded,
\[
n\log\left(1-\frac{P_n(t)}n\right)
=
%
-P_n(t)-\frac{P_n(t)^2}{2n}+O_t(n^{-2}).
\]
It follows that
\[
B_n(t)
=
%
e^{-P_n(t)}
\left(
1-\frac{P_n(t)^2}{2n}+O_t(n^{-2})
\right).
\]
Consequently,
\[
\boxed{
n\bigl(B_n(t)-\beta(t)\bigr)
\longrightarrow
-\frac12P(t)^2e^{-P(t)}.}
\tag{6.1}
\]

For example, take \(t=1/2\), corresponding to \(x=1/\sqrt2\). Then
\[
P(1/2)=1,\qquad \beta(1/2)=e^{-1},
\]
and
\[
n\left(f_n^{(b)}!\left(\frac1{\sqrt2}\right)
-\alpha\!\left(\frac1{\sqrt2}\right)\right)
\longrightarrow-\frac1{2e}.
\]
Therefore
\[
|f_n^{(b)}-\alpha|_\infty
\not=o(n^{-1}),
\]
so the \(O(n^{-1})\) rate is optimal.

\medskip\hrule\medskip

\subsubsection{6.4. Role of the domain}

The restriction to ([-1,1]) is essential. If \(|x|>1\), then
\[
x^2+x^4+\cdots+x^{2n}
\asymp_x x^{2n},
\]
so
\[
\left\lvert 
1-\frac{x^2+x^4+\cdots+x^{2n}}n
\right\rvert ^n\longrightarrow+\infty.
\]
Thus the second sequence has no analogous convergence on any interval containing a point with \(|x|>1\).

The endpoint values \(\alpha(\pm1)=0\) are forced by the interior limit and, more strongly, make \(\alpha\) flat and \(C^\infty\) at the endpoints.

\medskip\hrule\medskip

\subsection{Conclusion}

For both proposed approximation schemes, not only the functions but every fixed derivative converges uniformly on the full closed interval:
\[
\boxed{
\forall k\in\mathbb N,\quad
\lim_{n\to\infty}
\left\lvert \bigl(f_n^{(a)}\bigr)^{(k)}-\alpha^{(k)}\right\rVert_\infty=0,
}
\]
and
\[
\boxed{
\forall k\in\mathbb N,\quad
\lim_{n\to\infty}
\left\lvert \bigl(f_n^{(b)}\bigr)^{(k)}-\alpha^{(k)}\right\rvert _\infty=0.
}
\]

Hence the answer to both parts is affirmative.
